\chapter{Related Work}
\label{ch:related}

The work in this thesis sits at the intersection of three lines of research: controllable text generation, energy-based and sampling-based decoding, and the study of the geometry of language-model representations. This chapter reviews each in turn, with an emphasis on the specific claims and assumptions that the experiments later test, and it closes by drawing out the gap in the literature that the thesis addresses.

\section{Controllable Text Generation}
\label{sec:rel-ctg}

Approaches to controllable generation can be organized by how much they change the underlying model. At one extreme, \textbf{fine-tuning methods} adjust the model's parameters to internalize the desired behaviour, either by supervised training on curated data or, more commonly today, by reinforcement learning from human feedback \citep{ouyang2022training}. These methods are effective and now standard, but the control they provide is fixed at training time; introducing a new attribute requires new data and further training. \textbf{Prompting and conditioning methods} such as CTRL \citep{keskar2019ctrl} steer generation with control codes prepended to the input, which is inexpensive but limited to attributes anticipated during training.

At the other extreme are \textbf{decoding-time methods}, which leave the model frozen and intervene only during generation. These are the methods most relevant to this thesis, because they are the ones that promise inference-time flexibility. PPLM \citep{dathathri2020plug} perturbs the hidden states of a frozen model in the direction of an attribute classifier's gradient at each step. FUDGE \citep{yang2021fudge} and GeDi \citep{krause2021gedi} reweight the next-token distribution using a discriminator that predicts whether a prefix will lead to the desired attribute. These methods remain autoregressive: they modify the local next-token decision but retain the strictly left-to-right structure and therefore its inability to revise. The sampling-based methods discussed next are an attempt to escape that structure entirely.

\section{Energy-Based and Sampling-Based Decoding}
\label{sec:rel-energy}

A distinct and influential line of work casts decoding as sampling from an energy defined over whole sequences, as in \eqref{eq:intro-energy}, and thereby allows global constraints to be imposed on the completed text rather than token by token. COLD \citep{qin2022cold} performs energy-based constrained decoding by running Langevin dynamics on a relaxed representation of the sequence, combining a fluency energy from the language model with task-specific constraint energies. MuCoLa \citep{kumar2022gradient} samples in the continuous embedding space and projects back to tokens, and it is the closest published antecedent of the continuous sampler studied here. Related efforts include gradient-based inference for infilling and lexically constrained generation \citep{sha2020gradient} and a range of discrete samplers adapted from the energy-based-model literature, such as the gradient-informed discrete proposals of \citet{grathwohl2021oops} and the discrete Langevin proposals of \citet{zhang2022langevin} on which the discrete sampler here is based.

A recurring and rarely foregrounded feature of this literature is that the Metropolis--Hastings correction is frequently omitted. Methods are often described as sampling from the energy while in practice running a biased, noise-augmented optimizer, a distinction discussed in Section~\ref{sec:bg-mh}. Where the correction is included, the difficulty of computing the reverse proposal on a landscape defined by projection is seldom examined. One contribution of this thesis is to take the correction seriously, to implement it faithfully for both samplers, and to measure what happens when it is applied to a landscape that violates the smoothness it presupposes.

\section{Amortized Sampling and GFlowNets}
\label{sec:rel-gflownet}

Generative Flow Networks were introduced by \citet{bengio2021flow} as a method for sampling discrete compositional objects in proportion to a reward, with the explicit goal of producing diverse high-reward samples rather than the single mode that reward-maximizing reinforcement learning tends to find. The framework was formalized and extended in subsequent work \citep{bengio2023gflownet}, and a family of training objectives was developed, including trajectory balance \citep{malkin2022trajectory} and sub-trajectory balance \citep{madan2023learning}, which improve credit assignment along generative trajectories. \citet{hu2024amortizing} applied GFlowNets to language models, framing tasks such as infilling and reasoning as sampling from an intractable posterior defined by the frozen model, and reported that GFlowNet fine-tuning yields diverse samples that transfer better than those from reward-maximizing baselines. The GFlowNet component of this thesis builds directly on their formulation and codebase, using the same modified sub-trajectory objective and the same parameter-efficient adaptation of the base model \citep{hu2022lora}. The question posed here is narrower and more critical than theirs: not whether amortized sampling is useful in general, but whether it escapes the specific pathology of the frozen-likelihood energy that defeats the gradient-guided samplers.

\section{The Geometry of Language-Model Representations}
\label{sec:rel-geometry}

A final line of work concerns the geometry of the embedding spaces in which the continuous sampler operates, and it supplies part of the explanation for the results in Chapter~\ref{ch:results}. A well-documented property of contextual and static embeddings from Transformer language models is \textbf{anisotropy}: embeddings do not spread uniformly through the representation space but occupy a narrow cone, so that randomly chosen tokens have high average cosine similarity \citep{gao2019representation, ethayarajh2019contextual}. Anisotropy has direct consequences for any method that treats Euclidean or cosine distance in embedding space as a meaningful measure of semantic similarity, because it compresses the dynamic range of those distances. This thesis measures the anisotropy of the two models it studies and shows that it renders Euclidean distance an unreliable evaluation metric, which is one reason the experiments rely on a divergence-based measure of contextual fit instead.

A second and complementary body of work concerns the failure of high-probability text to be high-quality text. \citet{holtzman2020curious} showed that the most probable sequences under a neural language model are often degenerate, repetitive, and unnatural, and that maximization-based decoding such as beam search produces markedly worse text than sampling. This observation, which the present thesis reproduces on its own models and terms the \emph{likelihood trap}, is central to the interpretation of the results: it means that any procedure which succeeds in concentrating probability mass on the lowest-energy regions of the frozen-likelihood energy is thereby concentrating it on degenerate text.

\section{The Gap Addressed by This Thesis}
\label{sec:rel-gap}

The literature reviewed above establishes energy-based sampling as an attractive route to controllable generation and provides both the samplers and the amortized alternative that this thesis studies. What it does not provide is a direct, controlled test of the assumption on which the whole enterprise rests, namely that the gradient of a frozen autoregressive likelihood is a usable search direction on the discrete text manifold, together with a mechanistic account of what happens when that assumption fails. Prior negative observations tend to be reported in passing, as brittleness or as tuning difficulty, rather than isolated and explained. This thesis addresses that gap. It designs an ablation that separates the direction of the gradient from its magnitude and thereby tests the assumption cleanly; it constructs diagnostic experiments that identify the geometric, statistical, and analytic mechanisms behind the result; and it shows that the same mechanism survives amortization, thereby locating the cause in the training objective of the language model rather than in any particular sampling algorithm.
